\section{Decomposition theorem for semi-small maps (Xueqing Wen)}%
\label{sec:Decomposition theorem for semi-small maps (Xueqing Wen)}

\begin{quest}
    Let \(X\) be a (manifold, variety, analytic space, etc). How can we compute \(H^*(X, \C)\)?
\end{quest}

Unfortunately, this question is too hard, so we will replace it with a new question.

\begin{quest}
    How can we detect \(H^*(X, \C)\) using smaller spaces?
\end{quest}

\begin{thm}[K\"unneth formula]
    We have \(H^*(X \times Y, \C) = H^*(X, \C) \otimes H^*(Y, \C)\).
\end{thm}

Unfortunately, most spaces are not products, so we consider the following question.

\begin{quest}
    Let \(f \colon X \to Y\) be a morphism. Can we detect the cohomology of \(X\) using the cohomology of \(Y\) and the cohomology of the fibers of \(f\)?
\end{quest}

This question is still very difficult, but at least we can study it. We will begin with the case where \(f\) is a fiber bundle between compact manifolds. The answer is of course the Leray or Serre spectral sequences, but before we get there, we need some preliminaries.

\begin{defn}
    A \(\ul{\C}_Y\)-local system on \(Y\) is a locally constant sheaf \(\mc{L}\) of vector spaces on \(Y\).
\end{defn}

For example, if \(Y\) is connected, then local systems correspond to representations of \(\pi_1(Y)\). Also, note that if \(Y\) is an algebraic variety, we will need to work with its analytification.

\begin{exm}
    Let \(f\) be a fiber bundle with fiber \(F\) as above. Then \(R^i f_* \ul{\C}_X\) is a local system on \(Y\) with fiber \(H^i(F, \C)\).
\end{exm}

\begin{prop}
    We have
    \[ H^*(X, \C) = H^*(X, \ul{\C}_X) = H^*(\mr{pt}, R a_{X*} \ul{\C}_X). \]
\end{prop}

\begin{thm}[Leray/Serre spectral sequence]
    We have a spectral sequence
    \[ E_2^{p,q} = H^p(Y, R^q f_* \ul{\C}_X) \Rightarrow H^{p+q}(X, \C). \]
\end{thm}

We will now consider the setting where \(f \colon X \to Y\) is a smooth projective map between smooth projective varieties.

\begin{thm}[Deligne {\cite{Deligne1968}}]
    We have the non-canonical isomorphism
    \[ R f_* \ul{\C}_X = \bigoplus_{q \geq 0} R^q f_* \ul{\C}[-q]. \]
    At the level of cohomology, we have
    \[ H^*(X, \C) = \bigoplus_{q \geq 0} H^{*-q}(Y, R^q f_* \ul{\C}). \]
\end{thm}

This is a consequence of the hard Lefschetz theorem for projective varieties and also states that the Serre spectral sequence degenerates at the \(E_2\) page. In topology, this fails for the Hopf fibration.

We will now consider the setting where \(f \colon X \to Y\) is a proper map between algebraic varieties. Then we have a \textit{Whitney stratification}\footnote{We explicitly refuse to define this.} \(Y = \bigsqcup_{\lambda} Y_{\lambda}\) such that \(f^{-1}(Y_{\lambda}) \to Y_{\lambda}\) is a \(C^{\infty}\) fibration in possibly singular varieties.

\begin{defn}
    A \(\ul{\C}_Y\)-module \(\mc{F}\) on \(Y\) is called a \textit{constructible sheaf} if there exists a Whitney stratification \(Y = \bigsqcup Y_{\lambda}\) such that \(\mc{F}|_{Y_{\lambda}}\) is a local system for all \(\lambda\).
\end{defn}

With this definition, we may consider the derived category \(\mc{D}_c^b(Y)\) of constructible sheaves on \(Y\). This assignment admits a six functor formalism, where the functors are \(f_*, f^*, f_!, f^!, \otimes, \mc{H}\mr{om}\), and satisfies Verdier duality via the functor
\[ D_Y \colon \mc{D}_c^b(Y) \to \mc{D}_c^b(Y). \]

\begin{defn}
    A complex \(\mc{F}^{\bullet} \in \mc{D}_c^b(Y)\) is a \textit{perverse sheaf} if it satisfies the following conditions:
    \begin{enumerate}
        \item (Support condition) We have \(\dim \on{supp} \mc{H}^{-i}(\mc{F}^{\bullet}) \leq i\) for all \(i\);
        \item (cosupport condition) We have \(\dim \on{supp} \mc{H}^{-i}(D_Y \mc{F}^{\bullet}) \leq i\) for all \(i\).
    \end{enumerate}
\end{defn}

\begin{rmks}\leavevmode
    \begin{enumerate}
        \item Perverse sheaves are not sheaves, but are complexes in degrees \([-\dim Y, 0]\).
        \item They have two origins. One is intersection homology, and the other is the Riemann-Hilbert correspondence.
    \end{enumerate}
\end{rmks}

We will need to use the following fact: the category \(\ms{Perv}(Y)\) of perverse sheaves on \(Y\) is an abelian category and is moreover the heart of a \(t\)-structure on \(\mc{D}_c^b(Y)\). In particular, for any \(\mc{G} \in \mc{D}_c^b (Y)\), we can consider the perverse cohomology \(^p\mc{H}^i(\mc{G}) \in \ms{Perv}(Y)\).

\begin{thm}[Beilinson-Bernstein-Deligne-Gabber {\cite{BBD1982}}]
    Let \(f \colon X \to Y\) be a proper morphism between algebraic varieties. Assume that \(X\) is smooth and connected.\footnote{This is not necessary, but if we remove the smoothness assumption we need to consider the intersection complex.} Then
    \[ R f_* \ul{\C}_X [\dim X] = \bigoplus_{i} {^p\mc{H}^i}(R f_* \ul{\C}_X)[-i]. \]
    Moreover, each \(^p\mc{H}^i (R f_* \ul{\C}_X)\) is semi-simple (meaning each subobject is a direct summand).
\end{thm}

\begin{defn}
    Let \(f \colon X \to Y\) be a proper surjective morphism with \(X\) smooth and irreducible. Let \(Y = \bigsqcup Y_{\lambda}\) be a stratification such that \(f^{-1}(Y_{\lambda}) \to Y_{\lambda}\) is a \(C^{\infty}\)-fibration. We say that \(f\) is \textit{semi-small} if for all \(y_{\lambda} \in Y_{\lambda}\), we have \(\on{codim}_Y Y_{\lambda} \geq 2 \dim f^{-1} (y_{\lambda})\). If equality holds, we say that \(Y_{\lambda}\) is \textit{relevant}. If the only relevant stratum is the dense one, then we say that \(f\) is \textit{small}.
\end{defn}

\begin{thm}[{\cite{DeCataldoMigliorini2002}}]
    If \(f \colon X \to Y\) is semi-small, then \(R f_* \ul{\C}_X [\dim X]\) is a perverse sheaf on \(Y\).
\end{thm}

\begin{proof}
    For \(y_{\lambda} \in Y_{\lambda} \subseteq Y\), we see that
    \[ R f_* \ul{\C}_X |_{y_{\lambda}} = \bigoplus H^i (f^{-1}(y_{\lambda}), \C). \]
    Note that this vanishes if \(i > 2 \dim f^{-1}(y_{\lambda}) \leq \on{codim}_Y (Y_{\lambda})\). If we vary over all strata \(Y_{\lambda}\), we obtain the support condition. Because \(f\) is proper, we have \(R f_* = f_!\), and by using \(D_Y R f_* = f_! D_X\), we obtain the cosupport condition.
\end{proof}

\begin{exms}
    We now give some examples of semi-small maps.
    \begin{enumerate}
        \item Any generically finite proper map between surfaces is semi-small.
        \item Let \(Y = \ab\{xy = t^{\ell+1}\} \subseteq \C^3\). Then there is a crepant resolution \(f \colon \tilde{Y} \to Y\) of singularities, and the fiber over the singular point is a chain of \(\P^1\)s, so it is semi-small (see also the previous example).
    \end{enumerate}
\end{exms}

\begin{thm}[Kaledin {\cite{Kaledin2006}}]
    A projective birational map from a nonsingular holomorphic symplectic variety is semi-small.
\end{thm}

\begin{exm}
    For example, the Springer resolution and the Grothendieck-Springer resolution are semi-small. In addition, the Hilbert-Chow morphism for the Hilbert scheme of points on a surface is semi-small.
\end{exm}

Let \(f \colon X \to Y\) be semi-small and consider a stratification \(Y = \bigsqcup Y_{\lambda}\). If \(Y_{\lambda}\) is relevant, consider
\[ \mc{L}_{\lambda} = R^{d_{\lambda}} f_* \ul{\C}_{f^{-1}(Y_{\lambda})}, \]
where \(d_{\lambda} = 2 \dim f^{-1} (y_{\lambda})\). Then \(\mc{L}_{\lambda}\) is semi-simple. Here, note that \(\mc{L}_{|\lambda} = H^{d_{\lambda}}(f^{-1}(y_{\lambda}), \C)\). Now let \(\mr{IC}(\mc{L}_{\lambda})\) be the intermediate extension of \(\mc{L}_{\lambda}[\dim Y_{\lambda}]\) on \(\ol{Y}_{\lambda}\), which is a perverse sheaf.

\begin{thm}[{\cite{DeCataldoMigliorini2002}}]
    Let \(f \colon X \to Y = \bigsqcup Y_{\lambda}\) be semi-small. Then there is a canonical decomposition
    \[ R f_* \ul{\C}_X[\dim X] = \bigoplus_{Y_{\lambda}\text{ relevant}} \mr{IC}(\mc{L}_{\lambda}). \]
\end{thm}

\begin{exm}
    Let \(G\) be a reductive group with Lie algebra \(\g\) and choose a Borel \(B\). Then
    \[ G \times_B \b = \tilde{\g} \xrightarrow{\pi} \g \]
    is small. Therefore, we have
    \[ R \pi_* \ul{\C}_{\tilde{\g}} [\dim \g] = \mr{IC}(\mc{L}). \]
    Recall that \(\tilde{\g}^{\on{rs}} \to \g^{\on{rs}}\) is a \(W\)-torsor. Then for \(a \in \g^{\on{rs}}\), we have \(\mc{L}_{|a} \cong \C[W]\), and therefore \(\mc{L} = R^0 \pi_* \ul{\C}_{\tilde{\g}^{\on{rs}}}\). This implies that \(W\) acts on \(\mc{L}\), and this action can be extended to \(\mr{IC}(\mc{L})\) and also on \(R \pi_* \C_{\tilde{\g}}\), and in particular acts on the fiber \(R \pi_* \C_{\tilde{\g}}|_a\) for any \(a \in \g\).
\end{exm}

\begin{exm}
    Consider the Springer resolution \(\mu \colon T^*(\on{GL}_3/B) \to \mc{N} = \mathbb{O}_{(1^3)} \sqcup \mathbb{O}_{(2,1)} \sqcup \mathbb{O}_{(3)}\). This map is semi-small, but unfortunately all nilpotent orbits are relevant. Note that the fiber above the \((2,1)\) orbit is a union of two \(\P^1\)s glued at a node, the fiber above the \((3)\) orbit is a point, and the fiber above the \((1^3)\) orbit is \(\on{GL}_3/B\).

    We see that \(\mc{L}_{(3)}\) must be trivial because \(\mu\) is an isomorphism above this orbit. Then \(\mc{L}_{(2,1)} = R^2 \mu_* \ul{\C}\) is a rank \(2\) local system on \(\mathbb{O}_{(2,1)}\). Finally, \(\mc{L}_{(1^3)} = R^6 \mu_* \ul{\C}\) is a rank \(1\) local system on \(\mathbb{O}_{(1^3)} = \mr{pt}\). Therefore, we have
    \[ R \mu_* \ul{\C}_{T^*\on{GL}_3/B}[6] = \mr{IC}(\mc{L}_{(3)}) \oplus \mr{IC}(\mc{L}_{(2,1)}) \oplus \mr{IC}(\mc{L}_{(1^3)}). \]
    Restricting to \(b \in \mathbb{O}_{(3)}\), we obtain \(\mc{L}_{(3)} = H^0(\mr{pt})\). Restricting to \(a \in \mathbb{O}_{(2,1)}\), we obtain
    \[ H^0(X, \C) \oplus H^2(X, \C) = \mr{IC}(\mc{L}_{(3)})|_{(a)} \oplus \mr{IC}(\mc{L}_{(2,1)})|_{(a)}. \]
    Restricting to \(0 = \mathbb{O}_{(1^3)}\), we obtain
    \begin{align*}
        H^*(\on{GL}_3/B, \C) = \C[W] &= \mr{IC}(\mc{L}_{(3)})|_0 \oplus \mr{IC}(\mc{L}_{(2,1)})|_0 \oplus \mc{L}_{(1^3)}|_0 \\
        &= H^0(\on{GL}_3/B, \C) \oplus \ab(\text{a summand of dimension } 4) \oplus H^6(\on{GL}_3/B, \C).
    \end{align*}
    As representations of \(W = S_3\), these correspond to the trivial representation, two copies of the standard representation, and the sign representation. For the Springer-theoretic construction and these decompositions, see~\cite{BorhoMacPherson1983}.
\end{exm}

